
先讨论式~\eqref{eq:7-4-57}
中的符号 $*(i)$。
涉及这一符号的项见
引理~\ref{lem:7-4-73}；
也可参看图~\ref{fig:7-4-9}。
引理~\ref{lem:7-4-73}
右端第一项为
\begin{equation}
\begin{aligned}
&\bigl[
\widehat{\overline{\mathfrak m}}_1,\,
\mathfrak G\circ
\widehat{\mathfrak m}_{k,\beta}
\circ\mathfrak G\circ Y_i
\bigr]^{\mathrm{sup}}\\
&\qquad+
\bigl[
\widehat{\overline{\mathfrak m}}_1,\,
\mathfrak G\circ Y_i\circ
\mathfrak G\circ
\widehat{\mathfrak m}_{k,\beta}
\bigr]^{\mathrm{sup}},
\end{aligned}
\tag{7.4.80}\label{eq:7-4-80}
\end{equation}
并与其他算子复合。
式~\eqref{eq:7-4-79}
含有
\begin{equation}
-\mathfrak G\circ
\bigl(
\widehat{\mathfrak m}_{k,\beta}\circ Y_i
-(-1)^{\deg Y_i}
Y_i\circ
\widehat{\mathfrak m}_{k,\beta}
\bigr).
\tag{7.4.81.1}\label{eq:7-4-81-1}
\end{equation}
更准确地说，
它对应于
$\delta(\mathfrak G)
=-(\operatorname{id}-\Pi_H)$
中的 $-\operatorname{id}$。
这里
\[
\delta(\mathfrak G)
=\bigl[
\overline{\mathfrak m}_1,
\mathfrak G
\bigr]^{\mathrm{sup}}
=\overline{\mathfrak m}_1\circ\mathfrak G
+\mathfrak G\circ\overline{\mathfrak m}_1;
\]
见引理~\ref{lem:5-4-28}。
在式~\eqref{eq:7-4-80}
的每一项中，
这里的 $\mathfrak G$
是出现的两个
$\mathfrak G$ 中位于第二处的那个。

引理~\ref{lem:7-4-73}
右端第二项为
\[
\bigl[
\widehat{\overline{\mathfrak m}}_1,\,
\mathfrak G\circ Y_i
\bigr]^{\mathrm{sup}},
\]
并与其他算子复合。
它含有
\begin{equation}
-\bigl(
\mathfrak G\circ
\widehat{\overline{\mathfrak m}}_1\circ Y_i
-(-1)^{\deg Y_i}
\mathfrak G\circ Y_i\circ
\widehat{\overline{\mathfrak m}}_1
\bigr).
\tag{7.4.81.2}\label{eq:7-4-81-2}
\end{equation}
由定义，
式~\eqref{eq:7-4-81-1}
与~\eqref{eq:7-4-81-2}
之和为
\begin{equation}
-\mathfrak G\circ\delta(Y_i).
\tag{7.4.81.3}\label{eq:7-4-81-3}
\end{equation}
这里的 $\delta$
是式~\eqref{eq:7-4-27}
中 Hochschild 复形的上边界算子。
式~\eqref{eq:7-4-81-3}
的负号在定义 $*(i)$
的公式中给出 $+1$；
其余的
$\deg'Y_k$ 项来自
$\mathfrak G\circ Y_k$
与 $\overline{\mathfrak m}_1$
的换序。

\par\medskip\noindent\begin{minipage}{\textwidth}
\centering
\setlength{\unitlength}{0.75mm}
\begin{picture}(70,43)
\linethickness{0.55pt}
\put(32,39){\line(0,-1){34}}
\put(32,29){\circle*{1.8}}
\put(32,14){\circle*{1.8}}
\put(36,29){$\scriptstyle Y$}
\put(35,22){$\scriptstyle\operatorname{id}$}
\put(36,13){$\scriptstyle\mathfrak m_{1,\beta}$}
\end{picture}
\par\smallskip
{\renewcommand{\thefigure}{7.4.12}\refstepcounter{figure}
\normalfont\bfseries 图~\thefigure\label{fig:7-4-12}}
\end{minipage}\par\medskip

考虑图~\ref{fig:7-4-12}
中的树 $\Gamma$
作为上述讨论的例子。
这里
\[
Y\in
\operatorname{Hom}\bigl(
B_1(\overline C[1]),\overline C[1]
\bigr).
\]
有
\[
\mathfrak g(\Gamma)(x)
=\Pi_H\bigl(
\mathfrak m_{1,\beta}(
\mathfrak G(Y(x)))
\bigr).
\]
于是，以下展开中的最后一项
采用经校核的复合次序\translatornote{源书最后一项的括号写成 $\overline{\mathfrak m}_1\circ Y-(-1)^{\deg Y}\circ\overline{\mathfrak m}_1$，第二个复合号左侧漏了 $Y$；依同页式 (7.4.81.2) 补为 $Y\circ\overline{\mathfrak m}_1$。见疑误 A-0576。}：
\[
\begin{aligned}
&\overline{\mathfrak m}_1
\bigl(\mathfrak g(\Gamma)(x)\bigr)
-(-1)^{\deg Y}
\mathfrak g(\Gamma)
\bigl(\overline{\mathfrak m}_1(x)\bigr)\\
&\quad=
\sum_{\substack{
\beta_1+\beta_2=\beta\\
\beta_i\ne0
}}
\Pi_H\Bigl(
\bigl(
\mathfrak m_{1,\beta_1}\circ
\mathfrak m_{1,\beta_2}
\bigr)
\bigl(\mathfrak G(Y(x))\bigr)
\Bigr)\\
&\qquad+
\Pi_H\Bigl(
\mathfrak m_{1,\beta}
\bigl(
(\operatorname{id}-\Pi_H)
\bigl(\mathfrak G(Y(x))\bigr)
\bigr)
\Bigr)\\
&\qquad+
\Bigl(
\Pi_H\circ
\mathfrak m_{1,\beta}\circ
\mathfrak G\circ
\bigl(
\overline{\mathfrak m}_1\circ Y
-(-1)^{\deg Y}
Y\circ\overline{\mathfrak m}_1
\bigr)
\Bigr)(x).
\end{aligned}
\]
最后一项包含 $\delta Y$。
注意这里并没有出现
式~\eqref{eq:7-4-81-3}
中的负号：
在计算式~\eqref{eq:7-4-69}
时，整体上另有一个负号。
它源自邻接 $v_0$ 的外边
没有配置 $\mathfrak G$：
其他边都配置
$\mathfrak G$，
唯有此边配置 $\Pi_H$。
另见后面的式~\eqref{eq:7-4-83-3}，
那里也会出现这一负号。

接下来讨论 $*(i,j)$，
即式~\eqref{eq:7-4-57}
左端第二项的符号。
如前所述，
这一项来自 $Y$ 型边 $e$
所配置的算子
$\mathfrak G$ 的微分。
见图~7.4.13。
更准确地说，
它对应于
\[
\bigl[
\overline{\mathfrak m}_1,
\mathfrak G
\bigr]^{\mathrm{sup}}
=-(\operatorname{id}-\Pi_H)
\]
中的
$-\operatorname{id}$。
\par\medskip\noindent
\begin{minipage}{0.49\textwidth}
\centering
\setlength{\unitlength}{0.7mm}
\begin{picture}(88,64)
\linethickness{0.55pt}
\put(43,56){\line(0,-1){9}}
\multiput(43,58)(0,2){3}{\line(0,1){0.8}}
\put(43,47){\line(0,-1){22}}
\qbezier(43,47)(27,39)(11,31)
\qbezier(43,47)(36,39)(29,31)
\qbezier(43,47)(48,39)(52,31)
\qbezier(43,47)(58,39)(72,31)
\qbezier(43,25)(30,16)(16,7)
\qbezier(43,25)(40,16)(37,7)
\qbezier(43,25)(50,16)(58,7)
\qbezier(43,25)(59,16)(74,7)
\multiput(32,5)(-1,-2){3}{\line(0,-1){0.8}}
\multiput(59,5)(1,-2){3}{\line(0,-1){0.8}}
\put(46,49){$\scriptstyle Y_j$}
\put(46,35){$\scriptstyle\mathfrak G$}
\put(46,24){$\scriptstyle Y_i$}
\end{picture}
\par\smallskip
{\renewcommand{\thefigure}{7.4.13.1}\refstepcounter{figure}
\normalfont\bfseries 图~\thefigure\label{fig:7-4-13-1}}
\end{minipage}\hfill
\begin{minipage}{0.49\textwidth}
\centering
\setlength{\unitlength}{0.7mm}
\begin{picture}(88,64)
\linethickness{0.55pt}
\put(43,56){\line(0,-1){9}}
\multiput(43,58)(0,2){3}{\line(0,1){0.8}}
\put(43,47){\line(0,-1){22}}
\qbezier(43,47)(27,39)(11,31)
\qbezier(43,47)(36,39)(29,31)
\qbezier(43,47)(48,39)(52,31)
\qbezier(43,47)(58,39)(72,31)
\qbezier(43,25)(30,16)(16,7)
\qbezier(43,25)(40,16)(37,7)
\qbezier(43,25)(50,16)(58,7)
\qbezier(43,25)(59,16)(74,7)
\multiput(32,5)(-1,-2){3}{\line(0,-1){0.8}}
\multiput(59,5)(1,-2){3}{\line(0,-1){0.8}}
\put(46,49){$\scriptstyle Y_i$}
\put(46,35){$\scriptstyle\mathfrak G$}
\put(46,24){$\scriptstyle Y_j$}
\end{picture}
\par\smallskip
{\renewcommand{\thefigure}{7.4.13.2}\refstepcounter{figure}
\normalfont\bfseries 图~\thefigure\label{fig:7-4-13-2}}
\end{minipage}\par\medskip


示意地说，
图~\ref{fig:7-4-13-1}
对应于
\[
\bigl[
\overline{\mathfrak m}_1,\,
\mathfrak G\circ Y_i\circ
\mathfrak G\circ Y_j
\bigr]^{\mathrm{sup}}
\]
中的一项，即
\begin{equation}
(-1)^{\deg'Y_i+1}
\mathfrak G\circ Y_i\circ Y_j.
\tag{7.4.82.1}\label{eq:7-4-82-1}
\end{equation}
这里出现
$\deg Y_i$，
是因为
$\overline{\mathfrak m}_1$
与 $Y_i$ 换了次序；
指数中的 $+1$
来自
$\overline{\mathfrak m}_1$
与 $\mathfrak G$
的换序。
又因
\[
\overline{\mathfrak m}_1\circ\mathfrak G
+\mathfrak G\circ\overline{\mathfrak m}_1
=-(\operatorname{id}-\Pi_H),
\]
还要得到一个负号。

图~\ref{fig:7-4-13-2}
对应于
\[
\bigl[
\overline{\mathfrak m}_1,\,
\mathfrak G\circ Y_j\circ
\mathfrak G\circ Y_i
\bigr]^{\mathrm{sup}}
\]
中的一项，即
\begin{equation}
(-1)^{\deg'Y_j+1}
\mathfrak G\circ Y_j\circ Y_i.
\tag{7.4.82.2}\label{eq:7-4-82-2}
\end{equation}
注意图~\ref{fig:7-4-13-1}
与图~\ref{fig:7-4-13-2}
中的树，其
$*(\Gamma)$
（见式~\eqref{eq:7-4-61}）
彼此相差
$\deg'Y_i\,\deg'Y_j$。
综上，
\begin{equation}
\begin{aligned}
&\mathfrak G\circ
\Bigl(
(-1)^{\deg'Y_i+1}
Y_i\circ Y_j
+(-1)^{
\deg'Y_j+1+\deg'Y_i\deg'Y_j
}
Y_j\circ Y_i
\Bigr)\\
&\quad=
(-1)^{\deg'Y_i+1}
\mathfrak G\circ
\Bigl(
Y_i\circ Y_j
-(-1)^{\deg Y_i\deg Y_j}
Y_j\circ Y_i
\Bigr)\\
&\quad=
2(-1)^{\deg'Y_i}
\mathfrak G\circ[Y_i,Y_j].
\end{aligned}
\tag{7.4.82.3}\label{eq:7-4-82-3}
\end{equation}
注意式~\eqref{eq:7-4-28}
中有一个负号。
式~\eqref{eq:7-4-82-3}
的系数 $2$
与对称化算子
$\operatorname{Sym}$
定义中的 $1/\ell!$
相抵；见式~\eqref{eq:7-4-46}。
于是得到
$*(i,j)$ 中的
$\deg'Y_i$ 项。
涉及 $*(i,j)$
的其他项，
由前述换序过程给出。
